\section{讨论}
\subsection{比较的有效范围}
对比控制了初始候选、边界、动力学、扰动样本与评价预算，移除了局部精修对终端误差的混合影响。PSO在当前参数下更充分地利用群体搜索信息；GA的变异和SA的提议温度、步长并未针对该目标调优。因此结果应表述为实现和配置的比较，而非优化范式的优劣证明。多起点SA返回96条链中的最好访问点，不能与同预算之外的单链结果直接混用。

五组配对种子的样本量有限，本文报告全部点、中位数及四分位范围，不将通过数解释为真实飞行成功概率，也不据此作总体显著性声明。全部规划共用一个扰动实现，下一阶段需要独立环境样本及更大的种子集合检验稳定性。控制实验同样只使用一个独立于规划的实际扰动实现。

\subsection{数值精度与模型误差}
搜索值显著小于复核值并不意味着复核降低了轨迹质量，而是揭示了同一控制在不同积分近似下的差异。本文将搜索返回控制原样复核，不以高精度目标进行后续参数调整。细化差是局部数值敏感性检查，不能证明全部未采样轨迹的误差，也不能替代物理模型有效性验证。

指数大气、固定阻力系数、二维质点和静止地球限制了物理外推。当前未约束动压、热流、过载和姿态执行速率，故不对轨迹的工程适用性作结论。目标仅最小化终端误差，较少燃料或更长滑行只能作为所得解的伴随特征，不能解释为燃料最优性或霍曼转移最优性的证据。

\subsection{预测控制的局限}
时变线性模型沿参考建立，其有效性取决于相关动力学尺度和偏差，而不是仅由“位置偏差除以地心距离”的比例决定。地面软约束不能确保所有输入及质量硬约束永远可行；原始非线性运动也不必精确满足线性预测。理想全状态反馈、无传感器噪声及延迟的假设使实验较实际飞行更有利。本文未构造终端不变集或鲁棒约束收紧，不宣称严格稳定性、递归可行性或任意扰动保证。

后续可在相同评估协议下比较不同算法算子和参数敏感性，并引入燃料目标、工程路径约束、多环境样本、状态估计及鲁棒预测控制。此类扩展需与当前实验分别报告，以免同时改变多项条件而失去可解释性。
